% Drawing 154/02-17 -- folder 154, batch 02, archive page 23 (PDF page 24).
% Transcription: figure en tête de la colonne de gauche : un arc portant, dans l'ordre, les points $x$, $b$, $a$
% Kind: arc with three marked points. Method 2 (schematic redraw from the facsimile). Status: ok.
% Notes: Pen arc with three ticks, labelled x, b, a from left to right; the left end is doubled by a short stroke (arrow-like), drawn as a plain end. Redrawn 2026-11-12 (issue 128): to check again.
% Coordinates: 1 unit = 1 mm = 1/100 of the crop width (bbox in page fractions: [0.12, 0.085, 0.34, 0.14]).
% Redrawn by Claude (claude-fable-5-2), issue #128, 2026-11-12; not checked by a human; not compiled (no TeX in the session).
% Rights: reproduces part of the work -- behind RENDER_TRANSCRIPTS, never published by itself.
\documentclass[tikz,border=3pt]{standalone}
\usepackage{amsmath,amssymb}
\usetikzlibrary{patterns,arrows.meta}
\tikzset{pencil/.style={black!45}, struck/.style={black!35}, mpt/.style={circle,draw,inner sep=1.3pt}, pt/.style={circle,fill,inner sep=0.9pt}, opt/.style={circle,draw,fill=white,inner sep=1.1pt}}
\begin{document}
\begin{tikzpicture}[x=1mm,y=1mm,line cap=round,line join=round,font=\scriptsize]
\draw plot[smooth, tension=0.55] coordinates {(7,-24) (13,-19.5) (20,-16.5) (30,-14.8) (40,-14.2) (50,-14.6) (60,-16) (70,-18.5) (80,-22) (90,-26.5) (97,-31)};
\draw[very thick] (12.4,-18.6) -- (13.6,-20.6);
\draw[very thick] (29.6,-13.6) -- (30.4,-16);
\draw[very thick] (50,-13.4) -- (50,-15.8);
\node[inner sep=1pt] at (10.5,-14.5) {\small $x$};
\node[inner sep=1pt] at (29,-9) {\small $b$};
\node[inner sep=1pt] at (49,-10.5) {\small $a$};
\end{tikzpicture}
\end{document}
